\selectlanguage{ngerman}%
\frenchspacing%
\begin{ledgroupsized}[r]{120mm}%
\footnotesize%
\pstart%
\noindent\textbf{Überlieferung:}%
\pend%
\end{ledgroupsized}%
\begin{ledgroupsized}[r]{114mm}%
\footnotesize%
\pstart%
\parindent -6mm%
\makebox[6mm][l]{\textit{L}}%
Aufzeichnung:
LH~XXXVII\;4 Bl.\;100. 
Ein Zettel (10\:x\:12,5\;cm),
Ränder zum Teil beschnitten;
Papiererhaltungsmaßnahmen.
Zwei vollseitig beschriebene Seiten.
\pend%
\end{ledgroupsized}%
\vspace{5mm}%
\begin{ledgroup}%
\footnotesize%
\pstart%
\noindent%
\textbf{Datierungsgründe:}
Die Datierung der Aufzeichnung N.\;\ref{60255} beruht auf %
inhaltlichen Betrachtungen.
Siehe die ausführliche Begründung auf S.\;\pageref{60255_dat}.%
\pend%
\end{ledgroup}%
\selectlanguage{latin}%
\frenchspacing%
\vspace{8mm}%
\normalsize%
\pstart%
\noindent%
\lbrack%
\href{\Scan{LH_37_04_100r}}%
{100\,r\textsuperscript{o}}%
\rbrack%
\hfill%
Liquidum transparens%
\indexS{liquidum transparens}
\hfill%
\phantom{\lbrack100\,r\textsuperscript{o}\rbrack}%
\pend%
\vspace*{0.5em}%
\pstart%
\noindent%
Omne liquidum uniforme%
\indexS{liquidum uniforme}
in quo satis multi magnique sunt%
\indexS{pori in liquido}
\edtext{%
pori %
est transparens.%
}{%
\lemma{pori}%
\Bfootnote{%
\textit{(1)}~ad %
\textit{(2)}~\textlangle--\textrangle\ %
\textit{(3)}~lumi %
\textit{(4)}~motum vel %
\textit{(5)}~pressione\indexS{pressio} %
\textit{(6)}~vim 
\textit{(7)}~est transparens.%
~\textit{L}%
}}
Quo\-ni\-am enim liquidum est,
ideo a calore fusum est,%
\indexS{liquidum fusum a calore}%
\indexS{calor fundens}
hoc est,
partes ejus varie inter se agitantur,%
\indexS{partes liquidi agitatae}
in qua
\edtext{%
agitatione%
\indexS{agitatio partium liquidi}%
\indexS{agitatio celerrima}
cum celerrima sit,
et intra ictum oculi%
\indexS{ictus oculi}%
\indexS{oculus}
multa dispositionum millia
in eadem recta designabili
continga\lbrack n\rbrack t,
inter caeteras aliquoties%
\indexS{dispositio contingens}}{%
\lemma{agitatione}%
\Bfootnote{%
\textit{(1)}~saepe %
\textit{(2)}~cum celerrima %
\lbrack...\rbrack\ recta designabili %
\bars{contingat \textit{ändert Hrsg.}} inter caeteras aliquoties%
~\textit{L}%
}}
continget etiam illa,
ut pori recta inter se respondeant.%
\indexS{pori inter se respondentes}%
\indexS{pori in liquido}
Quin
\edtext{%
imo etiam cum non respondent pori,%
}{%
\lemma{imo}%
\Bfootnote{%
\textit{(1)}~etsi non resp\textlangle on\textrangle\ %
\textit{(2)}~etiam cum non respondent pori,%
~\textit{L}%
}}
tamen quando parum a se declinant,
ita ut cadant intra latitudinem%
\indexS{latitudo perexigua}
\edtext{%
perexiguam,
perinde est ac si sibi responderent.%
\indexS{pori inter se respondentes}%
\indexS{pori in liquido}
Et statum%
\indexS{status varians}
quo intra exiguum tempus%
\indexS{tempus exiguum}
et exiguum locum%
\indexS{locus exiguus}
variant,
a nobis
qui ista\lbrack s\rbrack\ exiguitates discernere non possunt,%
\indexS{exiguitas indiscernibilis nobis}
haberi pro uniformi aliquo statu%
\indexS{status uniformis medius}
inter varietates istas%
\indexS{varietas}
medio apparet%
\lbrack:\rbrack%
}{%
\lemma{perexiguam,}%
\Bfootnote{%
\textit{(1)}~ut et declina %
\textit{(2)}~perinde est %
\lbrack...\rbrack\ responderent. Et %
\bars{statum \textit{erg.}} %
\textit{(a)}~quae %
\textit{(b)}~quo intra exiguum tempus %
\textit{(aa)}~con %
\textit{(bb)}~variant %
\textit{(cc)}~et exiguum locum variant, %
\bars{ea \textit{streicht Hrsg.}} a nobis %
\textit{(aaa)}~hab %
\textit{(bbb)}~qui %
\bars{ista \textit{ändert Hrsg.}} exiguitates discernere non possunt, %
\bars{a nobis \textit{streicht Hrsg.}} haberi pro %
\lbrack...\rbrack\ medio apparet%
~\textit{L}%
}}
tum
\edtext{%
ratione loci%
\indexS{ratio loci}%
}{%
\lemma{ratione loci}%
\Bfootnote{%
\textit{erg.\;L}%
}}
ex coloribus%
\indexS{color mistus}
qui per exiguas partes misti novos producunt,%
\indexS{color novus productus}
\edtext{%
tum vero ratione temporis%
\indexS{ratio temporis}
\lbrack ex\rbrack\
eleganti experimento%
\indexS{experimentum elegans}%
}{%
\lemma{tum}%
\Bfootnote{%
\textit{(1)}~elegant %
\textit{(2)}~vero ratione temporis %
\bars{ex \textit{erg. Hrsg.}} eleganti experimento%
~\textit{L}%
}}
rotae dentatae%
\indexS{experimentum rotae dentatae}%
\indexS{rota dentata}
quae
si celeriter agitetur,
ambitus ejus dentatus apparet instar nebulae transparentis.%
\indexS{nebula transparens}
Hinc
\edtext{%
si
\lbrack%
\href{\Scan{LH_37_04_100v}}%
{100\,v\textsuperscript{o}}%
\rbrack\ %
liquidum ex globulis vel cylindris rigidis constet,%
\indexS{liquidum e globulis rigidis}%
\indexS{liquidum e cylindris rigidis}%
\indexS{globulus rigidus}%
\indexS{cylindrus rigidus}
vel etiam ex laminis%
}{%
\lemma{si}%
\Bfootnote{%
\textit{(1)}~lam %
\textit{(2)}~liquid %
\lbrack\href{\Scan{LH_37_04_100v}}{100\,v\textsuperscript{o}}\rbrack\ %
\textit{(3)}~liquidum ex %
\lbrack...\rbrack\ ex laminis%
~\textit{L}%
}}
vel tabulis varie agitatis,%
\indexS{liquidum e tabulis agitatis}%
\indexS{liquidum e laminis agitatis}%
\indexS{tabula agitata}%
\indexS{lamina agitata}
semper eadem
\edtext{%
fere ubique erit pororum constitutio%
\indexS{constitutio pororum in liquido}%
\indexS{pori in liquido}%
}{%
\lemma{fere}%
\Bfootnote{%
\textit{(1)}~ubique erit mensur %
\textit{(2)}~semper erit %
\textit{(3)}~ubique erit pororum constitutio%
~\textit{L}%
}}
sufficiens ad Aethera transmittendum.%
\indexS{aether transmittendus}
\pend%
\newpage%
\pstart%
Quin etsi corpus liquidum ponatur%
\indexS{corpus liquidum}
constare ex filamentis%
\indexS{liquidum e filamentis flexilibus}%
\indexS{filamentum flexile}
vel laminis flexilibus,%
\indexS{liquidum e laminis flexilibus}%
\indexS{lamina flexilis}
instar
\edtext{%
funium vel%
}{%
\lemma{funium}%
\Bfootnote{%
\textit{(1)}~et %
\textit{(2)}~vel%
~\textit{L}%
}}
velorum,%
\indexS{funis}%
\indexS{velum}
idem tamen
\edtext{%
evenit%
\lbrack:\rbrack\
\edlabel{LH_37_04_100_100v_srpnts-1}%
nam filamenta illa%
\indexS{filamentum flexile}%
}{%
\lemma{evenit,}%
\Bfootnote{%
\textit{(1)}~idem eni %
\textit{(2)}~nam filamenta illa%
~\textit{L}%
}}
celerrime huc illuc discurrent,%
\indexS{filamentum discurrens celerrime}
et
\edtext{%
ob\lbrack d\rbrack urante%
}{%
\lemma{oburante}%
\Bfootnote{%
\textit{L\;ändert Hrsg.}%
}}
motu%
\indexS{motus obdurans}
\edtext{%
simul
intorquebuntur varie%
\indexS{filamentum intortum}
atque incedent undanter,%
\indexS{filamentum incedens undanter}%
}{%
\lemma{simul}%
\Bfootnote{%
\textit{(1)}~unda %
\textit{(2)}~intorquebuntur varie atque incedent undanter,%
~\textit{L}%
}}
instar serpentum repentium%
\lbrack;\rbrack%
\indexS{serpens repens}%
\edlabel{LH_37_04_100_100v_srpnts-2}
ut videmus fila a calore aliquando intorqueri,%
\indexS{filum intortum a calore}%
\indexS{calor intorquens fila}
ita
\edtext{%
innumeros seminannulos ubique in liquido formabunt,%
\indexS{semiannuli per quos lumen transit in liquido}
per quos aeque facile transit lumen,%
\indexS{lumen transiens in liquido}%
}{%
\lemma{innumeros}%
\Bfootnote{%
\textit{(1)}~c\textlangle--\textrangle\ %
\textit{(2)}~formabunt %
\textit{(3)}~seminannulos %
\textit{(a)}~, quibus %
\textit{(b)}~\textlangle locu\textrangle s %
\textit{(c)}~ubique in liquido formabunt, %
\textit{(aa)}~ubique %
\textit{(bb)}~i %
\textit{(cc)}~per quos %
\lbrack...\rbrack\ transit lumen,%
~\textit{L}%
}}
ac nos
per arbores in quincuncem dispositas%
\indexS{arbores dispositae in quincuncem}
lumen aliquod motum,%
\indexS{lumen motum per arbores}
vel ipsi moti,
videre possumus
(alternis enim apparebit aut disparebit,%
\indexS{lumen apparens aut disparens alternis}
quod si celer sit motus noster vel luminis%
\lbrack,\rbrack\
nec discerni poterit,
sed videbitur perpetuo lucere).
Ut enim tunc nos movemur aut lumen,
ita res eodem rediret,
si arbores moveri intelligerentur%
\indexS{arbores dispositae in quincuncem}
uniformi ubique varietate.%
\indexS{varietas uniformis}
Si cessante calore%
\indexS{calor cessans}
fluidum deprehendatur et connectetur,%
\indexS{fluidum deprehensum}%
\indexS{fluidum connexum}
nihilo minus transparens manebit%
\indexS{fluidum transparens}
\edtext{%
quia statum transparentiae ultimae servat.%
\indexS{status transparentiae}%
\indexS{transparentia}%
}{%
\lemma{quia}%
\Bfootnote{%
\textit{(1)}~ultimae %
\textit{(2)}~statum transparentiae ultimae servat.%
~\textit{L}%
}}
\pend%
\begin{edfigure}[3.0em][0.0em]
\includegraphics[width=0.340\textwidth]%
{gesamttex/edit_VIII-4/images/LH_37_04_100_d_100v.pdf}
\edcaption[0.75em]{%
\lbrack\textit{Fig.\:1}\rbrack%
}%
\end{edfigure}%
\vspace{4.0em}%
\pstart%
\noindent%
\centering%
*     *     *
\pend%
\newpage%
\selectlanguage{german}%
\frenchspacing%
\footnotesize%
\pstart%
\noindent%
\textls{Ausführliche Datierungsgründe:}%
\pagelabel{60255_dat}
Die undatierte Aufzeichnung N.\;\ref{60255} \textit{Liquidum transparens} ist einer mechanistischen Erklärung für die Transparenz einer Flüssigkeit gewidmet.
Ihr Träger weist weder Spuren eines Wasserzeichens noch sonstige äußere Merkmale, die eine eindeutige Datierung begründen könnten.
Das kleine Papierformat und möglicherweise auch der Schreibduktus deuten aber auf eine Abfassung des Textes nach Leibnizens Aufenthalt in Paris\indexO{Paris} 1672 bis 1676.
\pend%
\pstart%
Inhaltlich weist N.\;\ref{60255} eine bemerkenswerte Verwandschaft mit der Aufzeichnung \textit{Cur aqua frigore intumescat} (N.\;\ref{60248} im vorliegenden Band) auf, in der sich Leibniz kritisch mit Descartes'\indexN{\textls{Descartes} (Cartesius, des Cartes), René 1596\textendash1650} Ansichten über das Verhalten des Wassers auseinandersetzt.
Insbesondere Leibnizens Vorstellung, dass eine Flüssigkeit aus \textit{filamenta} bestehe, die sich schnellstens hin und her bewegen und ineinander biegen würden \textit{instar serpentum repentium} (N.\;\ref{60255}, \refpassage{LH_37_04_100_100v_srpnts-1}{LH_37_04_100_100v_srpnts-2}), erinnert an die aalförmigen Teilchen, mit deren Hilfe Descartes\indexN{\textls{Descartes} (Cartesius, des Cartes), René 1596\textendash1650} das Schwellen durchfrierenden Wassers erklären wollte (siehe N.\:\ref{60248}).
Diese inhaltliche Verwandschaft lässt die Vermutung einer gemeinsamen Entstehung in den frühen Hannoveranver\indexO{Hannover} Jahren zu.
Diese Vermutung ist ferner mit der Feststellung vereinbar, dass auch eine Auseinandersetzung mit Descartes'\indexN{\textls{Descartes} (Cartesius, des Cartes), René 1596\textendash1650} \textit{Dioptrique}\cite{00038} ausschlaggebend für die Aufzeichnung N.\;\ref{60255} gewesen sein könnte.
Mit der \textit{Dioptrique}\cite{00038} war Leibniz seit den frühen 1670er Jahren vertraut (siehe die Datierungsgründe zu N.\;\ref{60248}, \pageref{60248_dat}).
\pend%
\pstart%
Eine Datierung der Aufzeichnung N.\;\ref{60255} \textit{Liquidum transparens} auf die Zeit vor 1677 %
ist nicht auzuschließen, aber unwahrscheinlich.
Ein begründeter Terminus ante quem der Entstehung von N.\;\ref{60255} lässt sich nach heutigem Wissensstand nicht angeben.
\pend%
\selectlanguage{latin}%
\frenchspacing%
\normalsize%
